% Part: first-order-logic
% Chapter: beyond
% Section: higher-order-logic

\documentclass[../../../include/open-logic-section]{subfiles}

\begin{document}

\olfileid[tr]{fol}{byd}{hol}

\olsection{Yüksek Dereceden Mantık}

Birinci dereceden mantıktan ikinci dereceden mantığa geçmek, biçimsel dilin
içinde birinci dereceden !!{domain}{}deki nesne kümelerinden söz etmemizi
sağladı. Neden orada duralım? Örneğin üçüncü dereceden mantık, nesne
kümelerinin kümelerini, hatta belki hem nesneleri hem de nesne kümelerini
içeren kümeleri ele almamızı sağlamalıdır. Dördüncü dereceden mantık da bu
tür nesnelerin kümelerinden söz etmemize olanak verir. Tahmin etmiş
olabileceğiniz gibi bu düşünce istenildiği kadar yinelenebilir.

Uygulamada yüksek dereceden mantık çoğu kez bağıntılar yerine fonksiyonlar
cinsinden !!{formula}{}leştirilir. (Doğal özdeşleştirmeler hesaba katıldığında bu fark
önemsizdir.) Bazı temel $A$, $B$, $C$,~\dots ``sıraları'' (bundan böyle bunlara
``tipler'' diyeceğiz) verildiğinde, şu koşulu koyarak yenilerini
oluşturabiliriz:
\begin{quote}
$\sigma$ ve $\tau$ sonlu tiplerse, $\sigma \to \tau$ da sonlu bir tiptir.
\end{quote}
Tipleri, !!{domain}{}imizde bulunmasını istediğimiz nesneleri sınıflandıran
sözdizimsel ``etiketler'' olarak düşünün; $\sigma \to \tau$, $\sigma$~tipindeki
nesneleri $\tau$~tipindeki nesnelere götüren fonksiyonlar olan nesneleri
betimler. Örneğin ``doğru'' ve ``yanlış'' doğruluk değerlerinden oluşan bir
$\Omega$ tipi ve doğal sayılardan oluşan bir $\Nat$ tipi isteyebiliriz. Bu
durumda $\Nat \to \Omega$ tipindeki nesneler tekli bağıntılar ya da $\Nat$'ın
alt kümeleri; $\Nat \to \Nat$ tipindeki nesneler doğal sayılardan doğal
sayılara fonksiyonlar; $(\Nat \to \Nat) \to \Nat$ tipindeki nesnelerse
fonksiyonları sayılara götüren yüksek tipli fonksiyonlar, yani
``fonksiyoneller'' olarak düşünülebilir.

İkinci dereceden mantıkta olduğu gibi yüksek dereceden mantık da ele almak
istediğimiz her nesne tipi için bir sıranın bulunduğu çok sıralı bir mantık
türü olarak düşünülebilir. Ancak yüksek tipli mantığın sözdizimini doğrudan
temelden tanımlamak genellikle daha açıktır. Örneğin bir sonlu tipler kümesini
tümevarımlı olarak şöyle tanımlayabiliriz:
\begin{enumerate}
\item $\Nat$ sonlu bir tiptir.
\item $\sigma$ ve $\tau$ sonlu tiplerse, $\sigma \to \tau$ da sonlu bir
  tiptir.
\item $\sigma$ ve $\tau$ sonlu tiplerse, $\sigma \times \tau$ da sonlu bir
  tiptir.
\end{enumerate}
Sezgisel olarak $\Nat$ doğal sayıların tipini, $\sigma \to \tau$ $\sigma$'dan
$\tau$'ya fonksiyonların tipini, $\sigma \times \tau$ ise biri $\sigma$'dan,
öteki $\tau$'dan olan nesne çiftlerinin tipini gösterir. Ardından terimler
kümesini tümevarımlı olarak şöyle tanımlayabiliriz:
\begin{enumerate}
\item Her $\sigma$ tipi için $x$, $y$, $z$, \dots değişkenlerinden oluşan,
  $\sigma$ tipinde bir dağarcık vardır.
\item $\Obj 0$, $\Nat$ tipinde bir terimdir.
\item $\Obj S$ (ardıl), $\Nat \to \Nat$ tipinde bir terimdir.
\item $s$, $\sigma$ tipinde bir terim ve $t$, $\Nat \to (\sigma \to
  \sigma)$ tipinde bir terimse, $\Obj{R}_{st}$, $\Nat \to \sigma$ tipinde
  bir terimdir.
\item $s$, $\tau \to \sigma$ tipinde ve $t$, $\tau$~tipinde bir terimse,
  $s(t)$, $\sigma$ tipinde bir terimdir.
\item $s$, $\sigma$~tipinde bir terim ve $x$, $\tau$~tipinde bir değişkense,
  $\lambd[x][s]$, $\tau \to \sigma$ tipinde bir terimdir.
\item $s$, $\sigma$~tipinde ve $t$, $\tau$~tipinde bir terimse,
  $\tuple{s, t}$, $\sigma \times \tau$ tipinde bir terimdir.
\item $s$, $\sigma \times \tau$~tipinde bir terimse, $p_1(s)$,
  $\sigma$~tipinde ve $p_2(s)$, $\tau$~tipinde bir terimdir.
\end{enumerate}
Sezgisel olarak $\Obj{R}_{st}$, özyinelemeyle
\begin{align*}
\Obj{R}_{st}(0) & = s \\
\Obj{R}_{st}(x+1) & = t(x, R_{st}(x)),
\end{align*}
biçiminde tanımlanan fonksiyonu; $\tuple{s, t}$, ilk bileşeni~$s$, ikinci
bileşeni~$t$ olan çifti; $p_1(s)$ ve~$p_2(s)$ ise $s$'nin birinci ve ikinci
elemanlarını (``izdüşümlerini'') gösterir. Son olarak $\lambd[x][s]$, $f$
fonksiyonunu gösterir; bu fonksiyon
\[
f(x) = s
\]
eşitliğiyle herhangi bir~$x$ için tanımlanır; bu değişken $\tau$ tipindedir.
Böylece (6).
madde, terimleri
kullanarak fonksiyon tanımlamamıza olanak veren bir kavrama biçimi sağlar.
!!^{formula}s, aynı tipteki terimler arasındaki $\eq[s][t]$ !!{identity}
önermelerinden, olağan önermesel bağlaçlardan ve yüksek tipli nicelemeden
kurulur. Sistemin aksiyomları da yukarıda tanımlanan terimleri yöneten temel
eşitliklerle, niceleyiciler ve !!{identity} içeren olağan mantık kurallarından
oluşturulabilir.

Sonlu tip sistemi bir doğruluk değerleri tipi $\Omega$ ile genişletilirse,
onun kullanımını yöneten aksiyomlar da eklenmelidir. Aslında yeterince ustaca
davranılırsa karmaşık !!{formula}s bütünüyle ortadan kaldırılıp yerlerine
$\Omega$ tipinde terimler konabilir!{} Kanıt sistemi de buna uygun biçimde
değiştirilebilir. Ortaya çıkan sistem özünde Alonzo Church'ün 1930'larda ortaya
koyduğu \emph{basit tipler kuramı}dır.

İkinci dereceden mantıkta olduğu gibi, yüksek tipli semantiğin de kullanılmak
istenebilecek farklı sürümleri vardır. Tam sürümde $\sigma \to \tau$ tipindeki
değişkenler, $\sigma$~tipindeki nesnelerden $\tau$~tipindeki nesnelere giden
\emph{bütün} fonksiyonlar kümesi üzerinde değer alır. Beklenebileceği üzere bu
semantik, tam ve etkin bir !!{derivation} sistemine izin vermeyecek kadar
güçlüdür. Ancak daha zayıf bir semantik ele alınabilir; burada !!a{structure},
$T_\tau$ eleman kümelerinden (her $\tau$ tipi için bir tane) ve uygulama,
izdüşüm vb. için uygun işlemlerden oluşur. Ayrıntılar doğru biçimde yürütülürse
yukarıda
betimlenen türden !!{derivation} sistemleri için tamlık teoremleri elde
edilebilir.

Yüksek tipli mantık çekicidir; çünkü matematiğin büyük bir bölümünü doğal
biçimde içine yerleştirebileceğimiz bir çerçeve sağlar: $\Nat$ ile başlayıp
gerçek sayılar, sürekli fonksiyonlar vb. tanımlanabilir. Tiplerin açık
``yapıcı'' yorumları bulunduğu için sezgici mantık bağlamında da özellikle
çekicidir. Nitekim bu yapıcı yorumları çok güzel açıklığa kavuşturan ve birçok
yönden klasik karşılıklarından daha ilginç olan yüksek tipli semantiğin yapıcı
sürümleri (klasik mantık yerine sezgici mantığa dayanarak) geliştirilebilir.

\end{document}
